% This file is generated from the corresponding Obsidian note.
% Source: Teoricos/GS - Teo14.md
\chapter{Extensión y caracterizaciones de suavidad}
\label{ch:teo-14}
\begin{bookexample}{Ejercicio abierto}
Sea $F:\mathbb R^2\to\mathbb R$ y $C=F^{-1}(0)$. $F(x,y)=x^2-y^3$. Se tiene $C=F^{-1}(0)$. ¿Se podria mostrar con la teoria vista hasta aca que $C$ NO puede ser subvariedad incrustada de $\mathbb R^2$? $C$ recibe el nombre de cuspide cubica. Si usamos este resultado: Suppose $M$ is a smooth manifold, $S\subseteq M$ is an embedded submanifold, and $p\in S$. As a subspace of $T_pM$, the tangent space $T_pS$ is characterized by $T_pS=\{v\in T_pM:vf=0\text{ whenever }f\in C^\infty(M)\text{ and }f|_S=0\}$. Es inmediata la solucion.

\end{bookexample}
\section{Extension de funciones suaves}
\begin{bookproposition}{Extension local}
Sea $(M,i)$ una inmersion en $N$ y sea $p\in M$ y sea $g:M\to\mathbb R$ funcion suave. Entonces, existen un abierto $U$ de $M$ en $p$, un abierto $V$ de $N$ en $i(p)$ con $i(U)\subseteq V$ (cuidado, $M$ no necesariamente tiene topologia heredada) y una funcion

\[
\widetilde g:V\subseteq N\to\mathbb R
\]
suave que extiende a $g$, es decir, $\widetilde g\circ i=g$ para todo $q\in U$. 
\bookimage{assets/pasted-image-20260505230207.png}{Pasted image 20260505230207}


\end{bookproposition}
\begin{bookproof}{}
\begin{enumerate}
\item Sea $m=\dim M$ y $n=\dim N$ si $m=n$ sale gratis (por que?)
\item Como $(M,i)$ es una inmersion, sabemos por la formula local de una inmersion que existen cartas $(U,\varphi)$ y $(V,\psi)$ centradas en $p$ e $i(p)$ tales que
\end{enumerate}
\[
\psi\circ i\circ\varphi^{-1}(x_1,\ldots,x_m)=(x_1,\ldots,x_m,0,\ldots,0)
\]
\begin{enumerate}[start=3]
\item Para poder proyectar adecuadamente, tenemos $\epsilon$ tal que $C_\epsilon^m(0)\subseteq\varphi(U)$ y $C_\epsilon^n(0)\subseteq\psi(V)$. Luego llamamos
\end{enumerate}
\[
U=\varphi^{-1}(C_\epsilon^m(0))\quad \text{y}\quad V=\psi^{-1}(C_\epsilon^n(0))
\]
\begin{enumerate}[start=4]
\item Como hemos visto en otras pruebas, tenemos $i(U)\subseteq V$ pues aplicando $\psi$ de ambos lados, esto equivale a
\end{enumerate}
\[
\psi\circ i\circ\varphi^{-1}(C_\epsilon^m(0))\subseteq C_\epsilon^n(0)
\]
y esto ultimo es cierto. 5. Si $(x_1,\ldots,x_m)\in C_\epsilon^m(0)$, entonces $|x_i|<\epsilon$, por lo tanto $(x_1,\ldots,x_m,0,\ldots,0)\in C_\epsilon^n(0)$. 6. Sea $\pi:\mathbb R^n\to\mathbb R^m$, $(y_1,\ldots,y_m,y_{m+1},\ldots,y_n)\mapsto(y_1,\ldots,y_m)$. 7. Nos animamos a proponer $\widetilde g:V\subseteq N\to\mathbb R$ dada por

\[
\widetilde g=g\circ\varphi^{-1}\circ\pi\circ\psi
\]
\begin{enumerate}[start=8]
\item Como
\end{enumerate}
\[
q\in V\xrightarrow{\psi}\psi(q)\in C_{\epsilon}^{n}(0) \xrightarrow{\pi}\pi(\psi(q))\in C_{\epsilon}^{m}(0) \xrightarrow{\varphi^{-1}}\varphi^{-1}(\pi(\psi(q)))\in U\xrightarrow{g}g\circ\varphi^{-1}\circ\pi\circ\psi(q)=\tilde{g}(q)
\]
, la funcion $\widetilde g$ claramente es funcion suave (por composicion de suaves) sobre $V$ . 9. Ademas si $q\in U$,

\[
\widetilde g\circ i(q)=\tilde g\circ i\circ \varphi\circ\varphi ^{-1}(q)=g\circ\varphi^{-1}\circ\pi\circ\psi\circ i\circ\varphi ^{-1}\circ\varphi(q)=g\circ\varphi ^{-1}\circ Id_{C_{\epsilon}^{m} (0)}\circ\varphi(q)=g(q)
\]
\bookqed
\end{bookproof}
\phantomsection\label{blk:f11484}
\begin{bookdefinition}{Suavidad sobre un subconjunto}
Sea $X$ una variedad suave y sea $A\subseteq X$. Una función $F:A\to \mathbb R^k$ se dice suave como función definida sobre el subconjunto $A\subseteq X$ si para todo $p\in A$, existe un abierto $O_p\subseteq X$ con $p\in O_p$ y existe una función suave $\widetilde F_p:O_p\to \mathbb R^k$ tal que

\[
\widetilde F_p|_{O_p\cap A}=F|_{O_p\cap A}.
\]
Es decir, $F$ es suave sobre $A$ si localmente alrededor de cada punto de $A$ se puede extender a una función suave definida en un abierto de la variedad ambiente $X$.

\end{bookdefinition}
\phantomsection\label{blk:c5f3d4}
\begin{bookproposition}{Extension de funciones suaves}
Sea $N$ una variedad suave, $M$ un subconjunto cerrado de $N$ y $F:A\to\mathbb{R}^k$ suave, en el sentido de \hyperref[blk:c5f3d4]{Teórico 14}. Entonces, para todo abierto $U$ que contenga a $M$, existe $\widetilde F:N\to\mathbb{R}^k$ función suave tal que

\[
\begin{cases}\widetilde F|_M=F,\\\operatorname{supp}\widetilde F\subseteq U.\end{cases}
\]
\end{bookproposition}
\phantomsection\label{blk:81237c}
\begin{bookproposition}{Extension Global}
Sea $(M,i)$ una incrustacion de $N$, donde $i$ es inclusion y $M$ es un subconjunto cerrado de $N$. Dada $g \in C^\infty(M)$, existe $\widetilde g \in C^\infty(N)$ que extiende a $g$, es decir $\widetilde g \circ i=g$.

\end{bookproposition}
\begin{bookproof}{}
\begin{enumerate}
\item Vamos a usar \hyperref[blk:81237c]{Teórico 14}
\item Para aplicarlo hay que ver que $g \in C^\infty(M)$ implica que $g$ es suave sobre $M$ al verla como subconjunto de $N$.
\item Y si!. Dado $p \in M$, por \hyperref[blk:f11484]{Teórico 14} existe un abierto $U_p$ de $p$ en $M$, un abierto $V_p$ de $p$ en $N$ y una funcion suave $\widetilde g_p:V_p \subset N \to \mathbb R$ tal que
\end{enumerate}
\[
g=\widetilde g_p \circ i:U_{p}\rightarrow \mathbb{R}
\]
sobre $U_p$. 4. Como $i$ es inclusion, esto dice que $U_p \subset V_p$ entonces $U_{p}\cap V_{p}=U_{p}$ 5. . Ademas, como $M$ esta incrustada en $N$, existe $W_p$ abierto de $N$ tal que

\[
U_p=M \cap W_p
\]
\begin{enumerate}[start=6]
\item Intersecando $W_{p}$ con $V_p$ tenemos que la funcion $\widetilde g_p|_{W_{p}\cap V_{p}}$ extiende a $g$ sobre $U_{p}=M\cap W_{p}\cap V_{p}$, un entorno de $p$ en $M$.
\item Con lo cual localmente osea un abierto $U_{p}$ al rededor de $p$ pudimos extender $g$ a una funcion $\tilde{g}_{p}$ suave sobre un abierto de $N$
\item Entonces $g$ es suave sobre $M$ como subconjunto de $N$.
\item Como $M$ es cerrado, aplica el resultado de extension global y se obtiene $\widetilde g \in C^\infty(N)$.
\end{enumerate}
\bookqed
\end{bookproof}
\section{Campos vectoriales}
\begin{bookremark}{}
Recordemos que $TM=\coprod_{p \in M}T_pM$ y que tenemos la proyeccion canonica $\pi:TM \to M$, $\pi(p,v)=p$.

\end{bookremark}
\begin{bookdefinition}{Seccion}
Sean $X$ e $Y$ conjuntos y sea $\pi:X \to Y$ una funcion. Una seccion de $\pi$ es una funcion $\sigma:Y \to X$ tal que $\pi \circ \sigma=\operatorname{Id}_Y$. Es decir, dado $q \in Y$, la seccion elige un elemento de la fibra $\pi^{-1}(q)$. (Al conjunto $\pi^{-1}(q)$ se lo llama la fibra de $\pi$ sobre $q$.)\\
En particular, si $\pi$ admite secciones, entonces $\pi$ es sobreyectiva. 
\bookimage{assets/pasted-image-20260509232539.png}{Pasted image 20260509232539}


\end{bookdefinition}
\begin{bookdefinition}{Campo vectorial}
Sea $M$ una variedad suave. Un campo vectorial sobre $M$ es una seccion de la proyeccion canonica $\pi:TM \to M$.

Dicho de otra forma, es una funcion $X:M \to TM$ tal que $\pi \circ X(q)=q$. Si escribimos $X(q)=(X_1(q),X_2(q))$, entonces

\[
X_1(q)=\pi\circ X(q)=q
\]
, y por lo tanto $X$ tiene la forma $X(q)=(q,v)$ con $v \in T_qM$. Al vector tangente asignado en $q$ se lo denota $X_q$, e identificamos $X(q)$ con $X_q$.

\end{bookdefinition}
\begin{bookdefinition}{Campo vectorial suave}
Un campo vectorial $X:M \to TM$ se dice suave si $X$ es una funcion suave, donde $TM$ tiene la estructura de variedad diferenciable vista anteriormente. Se denota al conjunto de campos suaves sobre $M$ por $\mathfrak X(M)$.

\end{bookdefinition}
\begin{bookremark}{}
Sea $X:M\rightarrow TM$ un campo vectorial. Sea $(U,\varphi=(x_1,\ldots,x_n))$ una carta suave de $M$. Para cualquier $p \in U$, el conjunto

\[
\left\{\frac{\partial}{\partial x_1}\bigg|_p,\ldots,\frac{\partial}{\partial x_n}\bigg|_p\right\}
\]
es una base de $T_pM$. Por lo tanto, si $X_p \in T_pM$, entonces

\[
X_p=\sum_{i=1}^n a_i(p)\frac{\partial}{\partial x_i}|_p
\]
Asi tenemos funciones $a_i:U \to \mathbb R$. Ademas se define

\[
\frac{\partial}{\partial x_{i}}:U\rightarrow  TU\qquad \text{dada por }\qquad \frac{\partial}{\partial x_{i}}(q)=\left(q,\frac{\partial}{\partial x_{i}}\bigg|_{q}\right)
\]
Luego a $a_{i}$ se llaman las funciones coordenadas de $X|_U$ con respecto al marco coordenado

\[
\left\{\frac{\partial}{\partial x_1},\ldots,\frac{\partial}{\partial x_n}\right\}
\]
Ademas, si $f \in C^\infty(U)$, definimos

\[
Xf:U \to \mathbb R
\]
por

\[
(Xf)(p)=X_p f
\]
\end{bookremark}
\begin{bookproposition}{Caracterizaciones de suavidad}
Sea $X:M \to TM$ un campo vectorial. Son equivalentes:

\begin{enumerate}
\item $X$ es suave.
\item Para toda carta suave $(U,\varphi=(x_1,\ldots,x_n))$ de $M$, las funciones $a_i:U \to \mathbb R$ dadas por
\end{enumerate}
\[
X_p=\sum_{i=1}^n a_i(p)\frac{\partial}{\partial x_i}|_p
\]
son suaves. 3. Para todo abierto $V$ de $M$ y toda funcion $f \in C^\infty(V)$, la funcion $Xf:V \to \mathbb R$, $(Xf)(q)=X_q f$, es suave. 4. Para toda funcion $f \in C^\infty(M)$, se cumple $Xf \in C^\infty(M)$.

\end{bookproposition}
\begin{bookproof}{}
\begin{itemize}
\item $1 \Rightarrow 2$.
\end{itemize}
\begin{enumerate}
\item Tenemos la carta $(U,\varphi=(x_{1},\ldots,x_{n}))$ la cual induce $(\widetilde{U}=\pi ^{-1}(U),\widetilde{\varphi})$ dada por
\end{enumerate}
\[
\widetilde\varphi(p,v)=(x_1(p),\ldots,x_n(p),v_1,\ldots,v_n)
\]
donde $v_i$ son las coordenadas del vector $v$ en la base $\{ \frac{\partial}{\partial x_{1}}|_{p},\ldots,\frac{\partial}{\partial x_{n}}|_{p} \}$ 2. Como $X$ es suave,

\[
\widetilde\varphi \circ X|_U:U \to \mathbb R^{2n}
\]
es suave. 3. Pero por la definición de las funciones $a_i$ tenemos que $X_p=\sum_{i=1}^n a_i(p)\frac{\partial}{\partial x_i}|_p$, osea las coordenadas de $X_p$ son $a_i(p)$. Por lo tanto

\[
\widetilde\varphi \circ X|_U(p)=\widetilde{\varphi}(p,X_{p})=(x_1(p),\ldots,x_n(p),a_1(p),\ldots,a_n(p))
\]
entonces cada funcion coordenada es suave. En particular, las $a_i$ son suaves.

\begin{itemize}
\item $2 \Rightarrow 3$.
\end{itemize}
\begin{enumerate}
\item Sea $V$ abierto de $M$ y sea $f \in C^\infty(V)$. Quiero ver $Xf:V\rightarrow\mathbb{R}$ es suave para esto alcanza con chequear suavidad con cartas suaves de $V$
\item Tomemos una carta $(U,\varphi=(x_1,\ldots,x_n))$ de $V$. Entonces para todo $p\in U$ tenemos
\end{enumerate}
\[
X_p=\sum_{i=1}^n a_i(p)\frac{\partial}{\partial x_i}|_p
\]
, con $a_i$ suaves. 3. Tenemos $Xf:U \to \mathbb R$ con

\[
(Xf)(p)=X_p f=\sum_{i=1}^n a_i(p)\frac{\partial }{\partial x_i}\bigg|_{p}f
\]
\begin{enumerate}[start=4]
\item Como las funciones coordenadas $a_{i}(p)$ son suaves basta ver que
\end{enumerate}
\[
\frac{\partial}{\partial x_{i}} f:U\subseteq M\rightarrow \mathbb{R}\quad\text{ dada por }\quad \frac{\partial}{\partial x_{i}}f(p)=\frac{\partial}{\partial x_{i}}\bigg|_{p}f
\]
son suaves (como funciones que reciben $p\in U$) 5. Pero lo son , por que

\[
\left(\frac{\partial}{\partial x_{i}}f\right)(p)=\frac{\partial}{\partial x_{i}}\bigg|_{p}f=\frac{\partial}{\partial r_{i}}\bigg|_{\varphi(p)}f\circ\varphi ^{-1}=\frac{\partial}{\partial r_{i}} f\circ\varphi ^{-1}(\varphi(p))
\]
osea

\[
\frac{\partial}{\partial x_{i}}f =\left(\frac{\partial}{\partial r_{i}}f\circ\varphi ^{-1}\right)\circ\varphi
\]
pero como $f\circ\varphi ^{-1}$ es suave entonces $\frac{\partial}{\partial r_{i}}f\circ\varphi ^{-1}$ es suave y como $\varphi$ es suave entonces $\frac{\partial}{\partial x_{i}}f$ es suave

\begin{itemize}
\item $3 \Rightarrow 1$.
\end{itemize}
\begin{enumerate}
\item Supongamos que para todo abierto $V\subseteq M$ y para toda función $f\in C^\infty(V)$, se tiene $Xf\in C^\infty(V)$. Queremos ver que $X:M\to TM$ es suave.
\item Para esto tomamos una carta suave $(U,\varphi=(x_1,\ldots,x_n))$ de $M$. Esta induce una carta suave del fibrado tangente
\end{enumerate}
\[
(\widetilde U=\pi^{-1}(U),\widetilde\varphi).
\]
\begin{enumerate}[start=3]
\item Entonces basta ver que
\end{enumerate}
\[
\widetilde\varphi\circ X|_U:U\to\mathbb R^{2n}
\]
es suave. (Esto es por que $p\in U$ y es directo ver que $X|_{U}(U)\subseteq \tilde{U}$ osea que despue es componer a derecha con $\varphi ^{-1}$ que es suave) 4. Para cada $p\in U$ tenemos

\[
\widetilde\varphi\circ X|_U(p)=\left(x_1(p),\ldots,x_n(p),a_1(p),\ldots,a_n(p)\right),
\]
donde $a_1(p),\ldots,a_n(p)$ son las coordenadas de $X_p$ en la base coordenada. 5. Las primeras $n$ componentes son suaves, porque $x_1,\ldots,x_n$ son las funciones coordenadas de una carta suave. 6. Por lo tanto, falta ver que

\[
a_i:U\to\mathbb R
\]
es suave para cada $i=1,\ldots,n$. 7. Para cada $p\in U$, tenemos

\[
X_p=\sum_{i=1}^n a_i(p)\frac{\partial}{\partial x_i}\bigg|_p.
\]
\begin{enumerate}[start=8]
\item Evaluamos en la función coordenada $x_i:U\to\mathbb R$. Entonces
\end{enumerate}
\[
X_px_i=a_i(p).
\]
\begin{enumerate}[start=9]
\item Por lo tanto,
\end{enumerate}
\[
a_i=Xx_i:U\to\mathbb R.
\]
\begin{enumerate}[start=10]
\item Pero $x_i\in C^\infty(U)$, entonces por hipótesis $Xx_i\in C^\infty(U)$. Luego $a_i$ es suave para cada $i$.
\item Así, todas las funciones componentes de $\widetilde\varphi\circ X|_U$ son suaves. Por lo tanto $\widetilde\varphi\circ X|_U$ es suave.
\item Como esto vale para cualquier punto $p$ del fibrado tangente, concluimos que $X:M\to TM$ es suave.
\end{enumerate}
\begin{itemize}
\item $3 \Rightarrow 4$ es inmediato tomando $V=M$.
\item $4 \Rightarrow 3$.
\end{itemize}
\begin{enumerate}
\item Queremos probar que si $V\subseteq M$ es abierto y $f\in C^\infty(V)$, entonces $Xf\in C^\infty(V)$.
\item Lo único que sabemos por hipótesis es que si $g\in C^\infty(M)$, entonces $Xg\in C^\infty(M)$.
\item Sea $p\in V$ arbitrario. Como $V$ es abierto, por \hyperref[blk:405b16]{Extensión de funciones suaves} existe un abierto $U\subseteq M$ tal que
\end{enumerate}
\[
p\in U \qquad \text{y} \qquad \overline{U}\subseteq V.
\]
y una funcion $\widetilde f\in C^\infty(M)$ tal que $\widetilde f|_{ U}=\left.f\right|_{ U}$ 4. Por hipótesis, como $\widetilde f\in C^\infty(M)$, entonces

\[
X\widetilde f\in C^\infty(M).
\]
\begin{enumerate}[start=5]
\item Además, como $\left.\widetilde f\right|_U=\left.f\right|_U$, entonces para todo $q\in U$ las funciones $\widetilde f$ y $f$ coinciden en un abierto alrededor de $q$. Por lo tanto,
\end{enumerate}
\[
X_q\left(\left.\widetilde f\right|_U\right)=X_q\left(\left.f\right|_U\right).
\]
por \hyperref[blk:b713ca]{Localidad de los vectores tangentes} 6. Como esto vale para todo $q$ tenemos que,

\[
(X\widetilde f)|_U=(Xf)|_U.
\]
\begin{enumerate}[start=7]
\item Como $X\widetilde f\in C^\infty(M)$, su restricción a $U$ es suave. Luego
\end{enumerate}
\[
(Xf)|_U\in C^\infty(U).
\]
\begin{enumerate}[start=8]
\item Como $p\in V$ era arbitrario, para cada punto de $V$ existe un abierto $U\subseteq V$ alrededor de ese punto tal que $(Xf)|_U$ es suave. Por lo tanto,
\end{enumerate}
\[
Xf\in C^\infty(V).
\]
notar que usamos \hyperref[blk:e32b2f]{Observacion}

\bookqed
\end{bookproof}
\phantomsection\label{blk:3accf8}
\begin{bookremark}{}
Como la suavidad de mapas entre variedades es una propiedad local en el dominio, para probar que un campo vectorial $X:M\to TM$ es suave alcanza probar que, para cada $p\in M$, existe un abierto $U\subseteq M$ con $p\in U$ tal que la restricción

\[
X|_U:U\to TM
\]
es suave.

Este primer paso no es especial de campos vectoriales: vale para cualquier mapa entre variedades. Lo específico del caso de campos vectoriales aparece cuando usamos que $X$ es una sección de $\pi:TM\to M$, es decir,

\[
\pi\circ X=\operatorname{id}_M.
\]
Entonces, si tomamos una carta suave $(U,\varphi=(x^1,\ldots,x^n))$ de $M$, se tiene automáticamente

\[
X(U)\subseteq \pi^{-1}(U),
\]
y podemos usar la carta inducida del fibrado tangente

\[
\widetilde\varphi:\pi^{-1}(U)\to\varphi(U)\times\mathbb R^n.
\]
En esa carta, si escribimos

\[
X_q=X^i(q)\frac{\partial}{\partial x^i}\bigg|_q,
\]
entonces la representación local de $X|_U$ es

\[
\widetilde\varphi\circ X\circ\varphi^{-1}(u)=\left(u^1,\ldots,u^n,X^1\circ\varphi^{-1}(u),\ldots,X^n\circ\varphi^{-1}(u)\right).
\]
Por lo tanto, probar que $X|_U$ es suave equivale a probar que las funciones coordenadas

\[
X^i:U\to\mathbb R
\]
son suaves.

\end{bookremark}
\phantomsection\label{blk:e32b2f}
